%% kapa: exam=1 round=29 year=2018 period=2 form=A subject=appraisal_law subject_ko=감정평가관계법규 seq=9 number=9
%% source: Q-Net 감정평가사 기출문제 (cst00309) artlSeq=5202808 file=extracted/1008/5202808/2018년 제29회 감정평가사 1차시험 기출문제/2018년 제29회 감정평가사 1차시험 기출문제/2교시 A형.pdf
%% conversion: kordoc 4.24.51 -> tools/gen1.py (auto); review=proofread-sonnet; flags=none
\begin{question}{9}
국토의 계획 및 이용에 관한 법령상 중앙도시계획위원회에 관한 설명으로 옳지 않은 것은?
\begin{choices}
  \item 국토교통부장관은 중앙도시계획위원회의 회의를 소집할 수 있다.
  \item 중앙도시계획위원회는 도시ㆍ군계획에 관한 조사ㆍ연구 업무를 수행할 수 있다.
  \item 중앙도시계획위원회의 회의는 재적위원 과반수의 출석으로 개의(開議)하고, 출석 위원 과반수의 찬성으로 의결한다.
  \item 중앙도시계획위원회의 위원장과 부위원장이 모두 부득이한 사유로 그 직무를 수행하지 못할 때에는 위원장이 미리 지명한 위원이 그 직무를 대행한다.
  \item 중앙도시계획위원회의 회의록은 심의 종결 후 3개월이 지난 후에는 공개요청이 있는 경우 이를 공개하여야 한다.
\end{choices}
\end{question}
\answer{5}
